
\documentclass{article}
\usepackage{colortbl}
\usepackage{makecell}
\usepackage{multirow}
\usepackage{supertabular}

\begin{document}

\newcounter{utterance}

\centering \large Interaction Transcript for game `codenames', experiment `abstract', episode 7 with reflexion{-}Qwen3.5{-}27B.
\vspace{24pt}

{ \footnotesize  \setcounter{utterance}{1}
\setlength{\tabcolsep}{0pt}
\begin{supertabular}{c@{$\;$}|p{.15\linewidth}@{}p{.15\linewidth}p{.15\linewidth}p{.15\linewidth}p{.15\linewidth}p{.15\linewidth}}
    \# & $\;$A & \multicolumn{4}{c}{Game Master} & $\;\:$B\\
    \hline

    \theutterance \stepcounter{utterance}  
    & & & \multicolumn{2}{p{0.3\linewidth}}{
        \cellcolor[rgb]{0.9,0.9,0.9}{
            \makecell[{{p{\linewidth}}}]{
                \texttt{\tiny{[GM$|$GM]}}
                {'hidden': {'team': ['expedience', 'docility', 'implication', 'romanticism', 'proliferation'], 'innocent': ['renown', 'antihuman', 'existentialism', 'refutation', 'treachery'], 'opponent': [], 'assassin': ['insubordination']}, 'revealed': {'team': {'team': [], 'innocent': [], 'opponent': [], 'assassin': []}, 'opponent': {'team': [], 'innocent': [], 'opponent': [], 'assassin': []}}}
            }
        }
    }
    & & \\ \\

    \theutterance \stepcounter{utterance}  
    & & \multicolumn{4}{p{0.6\linewidth}}{
        \cellcolor[rgb]{0.9,0.9,0.9}{
            \makecell[{{p{\linewidth}}}]{
                \texttt{\tiny{[P1$\langle$GM]}}
                \texttt{You have attempted this same task before and did not succeed. These are the plans you wrote after those attempts.} \\
\texttt{Trial \#1: New plan: My previous strategy failed because I used overly broad or ambiguous clues (e.g., "state," "growth," "soft") that allowed the partner to guess opponent words or distractors, ending my turn prematurely. Specifically, "state" led to "formality" (opponent), "growth" led to "furtherance" (opponent), and "soft" likely led to a distractor or opponent word given the loss. I also wasted turns by targeting only one word at a time when multiple words could have been grouped more safely.} \\
\\ 
\texttt{Next time, I will:} \\
\texttt{1.  **Prioritize Specificity:** Choose clues that are semantically tight to the target words but distinct from opponent/distractor words. For example, instead of "growth" (which overlaps with "furtherance"), I should have used "spread" or "multiplication" for "proliferation" alone, or grouped "proliferation" with "implication" under a clue like "consequence" if applicable, but only if the association is unique.} \\
\texttt{2.  **Avoid Overlap:** Before giving a clue, I will mentally check the opponent's list and distractors to ensure the clue does not strongly suggest any of those words. For instance, "soft" is risky because it could relate to "susceptibility" (opponent) or "docility" (target); a better clue for "docility" might be "obedience" or "meekness" to avoid the "soft" ambiguity.} \\
\texttt{3.  **Group Strategically:** Instead of single{-}word targets, I will look for strong semantic clusters. For the remaining words "expedience, docility, implication, romanticism, proliferation," I could try grouping "expedience" and "implication" under "pragmatism" or "logic" if they fit, or "romanticism" and "docility" under "idealism" or "gentleness" (checking for opponent overlap first).} \\
\texttt{4.  **Re{-}evaluate "Soft":** In the final turn, "soft" was too vague. I should have used "meek" for "docility" or "ideal" for "romanticism" to be more precise and avoid triggering the opponent's "susceptibility" or distractors.} \\
\\ 
\texttt{Specific new actions for the remaining words:} \\
\texttt{{-}   For "expedience": Use "pragmatism" or "convenience" (checking against "furtherance" or "foresight").} \\
\texttt{{-}   For "implication": Use "suggestion" or "inference".} \\
\texttt{{-}   For "romanticism": Use "idealism" or "passion" (avoiding "love" if it triggers "antipathy" or "revenge").} \\
\texttt{{-}   For "proliferation": Use "expansion" or "increase} \\
\\ 
\\ 
\texttt{Provide one single{-}word clue that relates to some of your team's words from the list below. You can choose to how many and to which words your clue relates to.} \\
\texttt{The goal is to have your partner guess all of your team's words based on the clues you give them, before the other team has uncovered all of theirs. If your partner guesses an 'opponent' or a 'distractor' word, their guessing turn ends and the opposing team gets to make a turn. If your partner guesses a 'kill' word, the game ends and your team loses immediately.} \\
\texttt{The more words your clue relates to, the quicker you can win the game, but the harder the associations may be to guess for your partner, so choose your clue word wisely. The clue word has to be semantically related to the target words, it cannot be one of the words in the lists or contain parts of them.} \\
\\ 
\texttt{Always give your single{-}word clue and your comma{-}separated list of related target words in the following format and make your answers as short as possible, never include any other text than is required in this form:} \\
\\ 
\texttt{CLUE: <WORD>} \\
\texttt{TARGETS: <TARGETS>} \\
\\ 
\texttt{Your team words are: expedience, docility, implication, romanticism, proliferation.} \\
\texttt{Your opponent's team words are: .} \\
\texttt{Distractor words are: renown, antihuman, existentialism, refutation, treachery.} \\
\texttt{The kill words are: insubordination.} \\
            }
        }
    }
    & \\ \\

    \theutterance \stepcounter{utterance}  
    & \multicolumn{4}{p{0.6\linewidth}}{
        \cellcolor[rgb]{0.9,0.9,0.9}{
            \makecell[{{p{\linewidth}}}]{
                \texttt{\tiny{[P1$\rangle$GM]}}
                \texttt{CLUE: consequence} \\
\texttt{TARGETS: implication, expedience} \\
            }
        }
    }
    & & \\ \\

    \theutterance \stepcounter{utterance}  
    & & & \multicolumn{2}{p{0.3\linewidth}}{
        \cellcolor[rgb]{0.9,0.9,0.9}{
            \makecell[{{p{\linewidth}}}]{
                \texttt{\tiny{[GM$|$GM]}}
                \texttt{consequence} \\
            }
        }
    }
    & & \\ \\

    \theutterance \stepcounter{utterance}  
    & & & \multicolumn{2}{p{0.3\linewidth}}{
        \cellcolor[rgb]{0.9,0.9,0.9}{
            \makecell[{{p{\linewidth}}}]{
                \texttt{\tiny{[GM$|$GM]}}
                ['implication', 'expedience']
            }
        }
    }
    & & \\ \\

    \theutterance \stepcounter{utterance}  
    & & & \multicolumn{2}{p{0.3\linewidth}}{
        \cellcolor[rgb]{0.9,0.9,0.9}{
            \makecell[{{p{\linewidth}}}]{
                \texttt{\tiny{[GM$|$GM]}}
                {'word': 'implication', 'assignment': 'team'}
            }
        }
    }
    & & \\ \\

    \theutterance \stepcounter{utterance}  
    & & & \multicolumn{2}{p{0.3\linewidth}}{
        \cellcolor[rgb]{0.9,0.9,0.9}{
            \makecell[{{p{\linewidth}}}]{
                \texttt{\tiny{[GM$|$GM]}}
                {'word': 'expedience', 'assignment': 'team'}
            }
        }
    }
    & & \\ \\

    \theutterance \stepcounter{utterance}  
    & & & \multicolumn{2}{p{0.3\linewidth}}{
        \cellcolor[rgb]{0.9,0.9,0.9}{
            \makecell[{{p{\linewidth}}}]{
                \texttt{\tiny{[GM$|$GM]}}
                \texttt{opponent has won} \\
            }
        }
    }
    & & \\ \\

    \theutterance \stepcounter{utterance}  
    & \multicolumn{4}{p{0.6\linewidth}}{
        \cellcolor[rgb]{0.95,0.95,0.95}{
            \makecell[{{p{\linewidth}}}]{
                \texttt{\tiny{[P2$\rightarrow$P1]}}
                \texttt{Post{-}episode Reflexion (not sent to GM):} \\
\texttt{New plan: My previous attempt failed because I prioritized grouping words over safety, selecting "consequence" which likely triggered the distractor "refutation" or "treachery" (both related to negative outcomes) or the kill word "insubordination" (a consequence of action), causing an immediate loss. I also left three words unaddressed. Next time, I will prioritize **safety and specificity** over grouping. I will select a clue that is semantically tight to a single target word to ensure the partner guesses it without ambiguity, avoiding any word that could loosely relate to "refutation," "treachery," or "insubordination." Specifically, for "docility," I will use "meekness" (distinct from distractors); for "romanticism," I will use "idealism" (distinct from "existentialism"); and for "proliferation," I will use "explosion" or "surge" (distinct from "renown"). I will only group words if the shared clue is exclusive to the targets and has zero semantic overlap with the opponent's list or distractors.} \\
            }
        }
    }
    & & \\ \\

\end{supertabular}
}

\end{document}
